% Generated draft: requires mathematical and source review.
\documentclass[11pt]{article}
\usepackage[margin=25mm]{geometry}
\usepackage{fontspec}
\setmainfont{Latin Modern Roman}
\usepackage{amsmath,amssymb,mathtools,unicode-math}
\setmathfont{Latin Modern Math}
\usepackage{enumitem,needspace,xurl,hyperref}
\setlength{\parindent}{0pt}
\setlength{\parskip}{5pt}
\setlength{\emergencystretch}{3em}
\newcommand{\N}{\mathbb{N}}
\newcommand{\Z}{\mathbb{Z}}
\newcommand{\Q}{\mathbb{Q}}
\newcommand{\R}{\mathbb{R}}
\newcommand{\C}{\mathbb{C}}
\newcommand{\F}{\mathbb{F}}
\newcommand{\A}{\mathbb{A}}
\newcommand{\PP}{\mathbb{P}}
\newcommand{\E}{\mathbb{E}}
\newcommand{\eps}{\varepsilon}
\newcommand{\dd}{\,\mathrm{d}}
\begin{document}
\section*{UnsolvedMath ID: 30006313}
\textbf{Source problem code:} catalogue entry
\subsection{Definitions and mathematical statement}
The supplied record contains two expository fragments and does not state a recoverable standalone open problem. The following retains their definitions and setup; it does not formulate the unspecified global duality conjecture.

For the base object denoted by \(\Sigma\) in the source, let \(\operatorname{Pic}_{\Sigma}\) denote the stack of line bundles on \(\Sigma\). The source introduces a stack \(\operatorname{Pic}^{X}_{\Sigma}\) classifying pairs \((L,\sigma)\), where \(L\) is a line bundle on \(\Sigma\) and \(\sigma\in H^{0}(\Sigma,L)\). No nonzero condition on \(\sigma\) is stated. Forgetting the section gives
\[
 p_X:\operatorname{Pic}^{X}_{\Sigma}\longrightarrow\operatorname{Pic}_{\Sigma},\qquad (L,\sigma)\longmapsto L.
\]
The source calls
\[
 (p_X)_!\overline{\Q}_{\ell}
\]
the \emph{period sheaf} and associates a function denoted \(\Theta_{\Phi}\) to it by the sheaf-to-functions correspondence. It says that this sheaf plays an important role in a global duality conjecture, but supplies no statement of that conjecture. The hypotheses on \(\Sigma\), the meaning of \(X\) and \(\Phi\), the coefficient prime \(\ell\), the sheaf-theoretic setting, and the Frobenius and normalization data for the correspondence are absent. Thus the displayed sheaf expression is retained as source notation with unresolved foundations. The later reference to a function \(\Theta_X\) and a zeta function does not specify whether \(\Theta_X\) and \(\Theta_{\Phi}\) coincide.

A subsequent fragment begins a separate discussion of the Gross--Prasad period. Let \(F\) be a global field and let \(V\) be a quadratic \(F\)-space with \(\dim_F V\geq 2\), not isomorphic to the hyperbolic plane. Write \(\langle 1\rangle\) for the one-dimensional quadratic space with underlying vector space \(F\) and form \(x\mapsto x^2\), and set
\[
 H=\operatorname{SO}(V),\qquad
 G=\operatorname{SO}(V)\times\operatorname{SO}(V\oplus\langle1\rangle).
\]
The indicated diagonal embedding is
\[
 \iota:H\longrightarrow G,\qquad h\longmapsto\bigl(h,h\oplus\operatorname{id}_{\langle1\rangle}\bigr).
\]
For every cuspidal tempered automorphic representation \(\pi\subset L^2([G])\) and every \(\varphi\in\pi\), the source introduces the period expression
\[
 P_H(\varphi)=\int_{[H]}\varphi(\iota(h))\,\mathrm{d}h.
\]
The conventional reading of the brackets is \([J]=J(F)\backslash J(\mathbb A_F)\) for \(J=G,H\), where \(\mathbb A_F\) is the adele ring of \(F\); this convention is not explicitly defined in the excerpt. The characteristic assumptions, precise quadratic-space conventions, Haar measure normalization, admissible realization of \(\varphi\), and convergence qualifications are not supplied. No convergence theorem, nonvanishing criterion, or conjectural identity is asserted here. The record supplies no mathematical question connecting this period setup to the preceding period-sheaf fragment.
\subsection{Short English statement}
The record defines a period sheaf by forgetting the section of a line bundle, then begins a separate Gross--Prasad period construction for special orthogonal groups. It supplies neither the global duality conjecture nor a standalone question. Both fragments are retained with their missing hypotheses and conventions explicitly unresolved; this is a statement-only draft requiring mathematical review.
\subsection{Sources}
\begin{itemize}
\item Supplied local UnsolvedMath registry record, permanent ID \(30006313\), exact source code \texttt{catalogue entry}, titled ``Period-Sheaf Expository Fragment,'' in the Oberwolfach Reports Open Problems collection. The record gives proposed year 2025 and identifies no proposer.
\item \emph{Arbeitsgemeinschaft: Relative Langlands Duality} (2025), Oberwolfach Reports, pp. 811--882, DOI \texttt{10.4171/owr/2025/18}: \url{https://doi.org/10.4171/owr/2025/18}. The extracted text includes the page marker 835 and the subsequent headings ``3. The Gross--Prasad period'' and ``3.1. The period integral.''
\item The excerpt refers to \S 1.4 of Sam Gunningham's article in these proceedings for the period sheaf and global duality, and to Shilin Lai's article for the relation between \(\Theta_X\) and the zeta function in the function-field setting. Separate article titles and bibliographic details are not supplied.
\item Supplied literature source: ``Arinkin et al., Geometric Langlands theory with nilpotent singular support,'' \url{https://arxiv.org/abs/2010.01906}. This attribution and link are retained as supplied; their bibliographic correspondence has not been independently verified.
\item Literature assessment dated \texttt{2026-08-22}: supplied dated curation, not independently verified. It labels the record open, describes it as an expository fragment rather than a resolved standalone theorem, and says the associated general global-duality framework remains conjectural. Its listed sources are the Oberwolfach report and the preceding arXiv link. The supplied background requests independent verification of time-sensitive status and cited claims.
\item The supplied statement-verification note reports that the official report confirms an extraction joining a period-sheaf definition to the setup for a different period problem, with no separate open problem. The registry marks the statement \texttt{unrecoverable}, provides no clean statement, and retains the original for provenance. This verification claim is reported from the record and was not independently checked here.
\end{itemize}
\end{document}
